% Part I -- Introduction and research framing (target: about 10 pages)
\clearpage
\part{Introduction and Research Framing}
\label{part:intro}
\partbudget{10}{opening argument, research questions, scope and method drafted; contributions finalised}

\section{Introduction}
\label{sec:intro}

\subsection{Motivation: why monolithic scaling is becoming difficult}
\label{sec:motivation}

For most of the history of the microprocessor, a new product generation could be built by placing more transistors on a single piece of silicon. The course textbook summarises this trend as Moore's law: the number of transistors on a chip has roughly doubled every one to two years since the 1960s \cite[sec.~12.1.1]{sarangi2025basic}.\index{Moore's law}

\subsubsection*{How processors used to scale}
The textbook describes three phases in how designers spent those transistors \cite[sec.~12.1.2]{sarangi2025basic}. First, a smaller transistor was faster and more power-efficient, so each generation built a bigger processor: more complex units, larger caches, wider issue and more functional units. Second, until about 2002, pipelines also grew deeper and clock frequencies rose with them. Third, after about 2002, power and temperature became the limit: processors began to exceed 100\,W and on-chip temperatures 100\,\textdegree C, which effectively ended the growth in core complexity and clock frequency. Designers instead began to place more cores of an unchanged design on each chip, so that the number of cores could double with each process generation \cite[sec.~12.1.2]{sarangi2025basic}.

Over the same period the smallest feature that could be manufactured shrank by a factor of about $\sqrt{2}$ every two years, from 130\,nm in 2001 to 22\,nm in 2011, doubling the number of transistors per unit area each time \cite[sec.~12.1.1, Table~12.1]{sarangi2025basic}. Writing around 2012, the textbook expected processors with 32 or 64 cores within five to ten years \cite[sec.~12.1.2]{sarangi2025basic}. The prediction came true almost exactly: AMD's second-generation EPYC offered 64 cores in 2019. But it did not arrive in the way the multicore era had assumed. Those 64 cores were not on one larger die but on eight small ones, joined in one package to a separate I/O die \cite{naffziger2020chiplet}. Explaining why is the starting point of this paper.

Three pressures now make it harder to keep growing a single \gls{monolithic-die}:

\begin{itemize}
  \item \textbf{Die size and yield.} A lithography system can only pattern a die up to the \gls{reticle-limit}, and the probability that a die contains a fatal defect rises with its area. AMD states that the 64-core server configuration of its second-generation EPYC processors would lie well beyond the maximum reticle size if built as one chip \cite{naffziger2020chiplet}. \Cref{sec:yield-cost} quantifies the yield effect with a simple model.
  \item \textbf{Uneven scaling of circuit types.} Logic, \gls{sram} and analog/I-O circuits do not shrink equally when moving to a new process node. Naffziger \emph{et al.} note that analog circuitry shrinks far less than digital logic and SRAM, so memory interfaces and SerDes gain little from an advanced node \cite{naffziger2020chiplet}. More recent AMD material indicates that SRAM itself now scales more slowly than logic (\cref{sec:why-cache-capacity}).
  \item \textbf{Power and temperature.} Dynamic power grows with switched capacitance, the square of the supply voltage and clock frequency, while leakage power is consumed even when circuits are idle and itself rises with temperature \cite[sec.~10.10]{sarangi2025basic}. Every additional square millimetre of active silicon must be powered and cooled within a fixed package budget.
\end{itemize}

These pressures do not stop progress, but they change its form (\cref{fig:zen-timeline}). Instead of one ever-larger die, designers increasingly assemble a processor from several dies inside one package, and increasingly use the vertical dimension as well as the horizontal one. This paper studies that shift through two concrete techniques and one principal case study.

\input{"../initial-draft-2/figures/zen-timeline.tex"}

\subsection{A thought experiment: one die or eight?}
\label{sec:thought-experiment}
Consider two ways to build a 64-core processor. Processor A places all 64 cores on one large die. Processor B places them on eight smaller dies of eight cores each, connected inside one package to a separate I/O die. Both offer 64 cores. Are they equally expensive to make? Do they communicate equally quickly? Can they use the same manufacturing process?

Using AMD's Zen~2 dimensions, with the compute dies at 74\,mm$^2$ each and a monolithic equivalent of $8 \times 74 = 592$\,mm$^2$ for the compute area alone, the later parts of this paper answer these questions (\cref{tab:thought-experiment}).

\begin{table}[htbp]
  \centering
  \caption{One die or eight: the 64-core thought experiment, answered from later sections (model results are illustrative)}
  \label{tab:thought-experiment}
  \small
  \begin{tabularx}{\textwidth}{@{}>{\raggedright\arraybackslash}p{3.0cm}>{\raggedright\arraybackslash}X>{\raggedright\arraybackslash}X>{\raggedright\arraybackslash}p{1.7cm}@{}}
    \toprule
    \textbf{Question} & \textbf{A: one 64-core die} & \textbf{B: eight 8-core dies + I/O die} & \textbf{Where} \\
    \midrule
    Can it be built? & For 64 cores with I/O, beyond the maximum reticle size \cite{naffziger2020chiplet} & Yes; each die is small & \cref{sec:monolithic-limits} \\
    Modelled yield ($D_0 = 0.2$\,cm$^{-2}$) & 38\% & 87\% per compute die & \cref{sec:yield-cost} \\
    64-core sets per 300\,mm wafer & about 35 & about 95 & \cref{sec:dies-per-wafer} \\
    AMD's cost estimate & roughly 2$\times$ at 32 and 48 cores & 1$\times$ (reference) & \cref{fig:chiplet-cost-chart} \\
    Manufacturing process & everything on the newest node & cores on 7\,nm, I/O on cheaper 12\,nm & \cref{sec:ccd-iod} \\
    Communication & on-die network only & every memory access crosses a die-to-die link (under 9\,ns round trip) & \cref{sec:d2d} \\
    Memory latency & not available for comparison & Rome's average about 0.81$\times$ that of its multi-die predecessor & \cref{tab:numa-latency} \\
    Cache per core & one shared design choice & each die has its own L3; can be enlarged per die by stacking & Part~\ref{part:stacked-cache} \\
    \bottomrule
  \end{tabularx}
\end{table}

The answers are not all in favour of B. It is far cheaper and the only option that can be built at all, but it pays for that with a die-to-die hop on every memory access and with caches that are private to each die. The rest of the paper examines both sides of this trade, and how 3D stacking changes the cache side of it.

\subsection{Two linked scaling problems}
\label{sec:two-problems}

The paper is organised around two related problems (\cref{fig:research-framework}).

\textbf{Problem 1: scaling processor resources beyond a single die.} If a product needs more cores, memory channels and I/O than one economical die can hold, the processor can be partitioned into \glspl{chiplet}. AMD's Zen~2-based products separate eight-core compute dies (CCDs) built on a 7\,nm process from an I/O die (IOD) built on a mature 12\,nm process \cite{naffziger2020chiplet,suggs2019zen2}.

\textbf{Problem 2: scaling cache capacity without spending planar die area.} A larger \gls{llc} can reduce costly trips to main memory, but on a planar die every extra megabyte of SRAM enlarges the die, lowers yield and lengthens wires. AMD's \gls{v-cache} instead stacks an additional SRAM die on top of the compute die, expanding the Zen~3 L3 from 32\,MB to 96\,MB per CCD \cite{wuu2022vcache,amd2022epyc7003}.

The problems are linked: the chiplet organisation makes the CCD a small, self-contained compute die with its own L3, and that die is the natural host for a stacked cache. The paper therefore follows the argument in four steps:

\begin{enumerate}
  \item Monolithic CPUs face die-size, yield and scaling limits; chiplets address compute scaling \cite{naffziger2020chiplet}.
  \item Chiplets alone leave the cost of large on-die caches unsolved: AMD's ISSCC~2022 presentation names SRAM scaling and cost as the barriers to large on-die caches, in cost, product flexibility and latency, and argues that 2.5D chiplets address only part of the problem \cite[slides~6--7]{wuu2022vcache}.
  \item Vertical stacking addresses those cache-specific limits \cite{wuu2022vcache,swaminathanAdvancedPackagingEnabling2021}.
  \item Evaluation evidence shows where the stacked cache improves performance, and which thermal, packaging and cost trade-offs it introduces (Part~\ref{part:case-study}).
\end{enumerate}

\input{"../initial-draft-2/figures/research-framework.tex"}

\subsection{Research questions and objectives}
\label{sec:rq}

The central research question is:

\begin{quote}
\itshape How do chiplet-based organisation and 3D-stacked cache help modern processors address limits in performance, scalability, memory latency, power and die-level integration, and what new trade-offs do they introduce?
\end{quote}

It is divided into four sub-questions, the second with a supporting question:

\begin{description}[leftmargin=1.2cm,style=nextline]
  \item[RQ1] Which architectural and manufacturing limits make a single large die an increasingly costly way to scale a processor?
  \item[RQ2] How does partitioning a processor into compute and I/O dies change memory latency, bandwidth, power and product flexibility?
  \item[RQ2a] How does AMD divide processor functionality between the CCDs and the IOD, and why does it connect them through the organic package substrate rather than a 2.5D or 3D interconnect?
  \item[RQ3] Why does last-level-cache capacity matter, and what does vertical stacking change compared with expanding the cache in the plane of the die?
  \item[RQ4] For which workloads does a larger stacked L3 produce measurable benefit, and how can that benefit be separated from differences in frequency, core count and memory configuration?
\end{description}

For every technique the paper applies the same three-part test: \emph{which bottleneck does it address, what benefit does it provide, and what new bottleneck, cost or trade-off does it create?} \Cref{tab:rq-map} shows where each question is answered.

\begin{table}[htbp]
  \centering
  \caption{Where the research questions are answered}
  \label{tab:rq-map}
  \small
  \begin{tabularx}{\textwidth}{@{}lX>{\raggedright\arraybackslash}p{4.2cm}@{}}
    \toprule
    \textbf{Question} & \textbf{Main evidence} & \textbf{Sections} \\
    \midrule
    RQ1 & Reticle limit, yield and dies-per-wafer models, uneven scaling of circuit types & Part~\ref{part:chiplets} (\cref{sec:monolithic-limits,sec:yield-cost,sec:dies-per-wafer}) \\
    RQ2, RQ2a & AMD's chiplet disclosures, die-to-die link figures, NUMA latencies, packaging options & Part~\ref{part:chiplets} (\cref{sec:ccd-iod,sec:d2d,sec:packaging-options}) \\
    RQ3 & Planar limits, hybrid bonding, the stacked-cache interface, thermal and manufacturing evidence & Part~\ref{part:stacked-cache} \\
    RQ4 & AMD-reported and independent measurements, graded for confounding factors & Part~\ref{part:case-study} \\
    All & Comparison, interactions and conclusions & Part~\ref{part:synthesis} \\
    \bottomrule
  \end{tabularx}
\end{table}

\subsection{Scope}
\label{sec:scope}

AMD Zen processors are the principal case study because AMD has published unusually detailed technical descriptions of both its chiplet organisation and its stacked cache \cite{naffziger2020chiplet,suggs2019zen2,amd2022epyc7003,swaminathanAdvancedPackagingEnabling2021}. They are, however, a case study and not the whole subject: the architectural principles apply to chiplet and die-stacked designs in general.

The following are deliberately outside scope:

\begin{itemize}
  \item fabrication chemistry and process details beyond what is needed to explain an architectural decision;
  \item reverse engineering of undisclosed Zen internals;
  \item a catalogue of product specifications;
  \item vendor contests. Comparisons are always tied to a stated workload, metric and time period, and the paper makes no general claim that one vendor ``beats'' another.
\end{itemize}

\subsection{Methodology}
\label{sec:method-overview}

The study combines three kinds of evidence:

\begin{enumerate}
  \item \textbf{Textbook foundations.} Chapters 9--12 of Sarangi's \emph{Basic Computer Architecture} supply the models used throughout: the performance equation, pipelining limits, power and temperature, the memory hierarchy and average memory access time, and multicore coherence and interconnects \cite{sarangi2025basic}.
  \item \textbf{Primary technical disclosures.} Peer-reviewed conference papers and vendor technical presentations describe what was built and why. Vendor presentations are attributed to the vendor and treated as engineering claims rather than independent measurements.
  \item \textbf{Published research and benchmarks.} Part~\ref{part:case-study} analyses published measurements of chiplet and stacked-cache processors. The paper does not perform new hardware measurements; its quantitative contribution lies in organising, normalising and critically comparing existing results (\cref{sec:evidence-method}).
\end{enumerate}

Two practices run through the analysis. First, \emph{worked examples} apply the textbook's models to published numbers: yield and cost per wafer, connection density, average memory access time, the stacked-cache trade-off and clock correction. Each states its inputs, and any assumed value is labelled as such. Second, sources are graded. Vendor presentations are used for what was built; independent measurements, preferably with fixed clocks or performance counters, are used for what it achieves; and where a value had to be read from a published chart rather than a table, the text says so and gives the reading accuracy.

\subsection{Contributions}
\label{sec:contributions}

This paper contributes no new measurements; its contributions are analytical:
\begin{itemize}
  \item a single argument linking chiplet partitioning and 3D-stacked cache as responses to two related scaling limits, die cost and last-level cache capacity;
  \item an explanation of both techniques in terms of the textbook's performance, power, thermal and memory-hierarchy models;
  \item a structured comparison of planar, side-by-side and vertically stacked cache expansion;
  \item a synthesis of AMD-reported and independent performance evidence for stacked cache across three Zen generations, graded by which confounding factors each comparison controls (\cref{tab:confound-grid});
  \item an account of how AMD's move to an inverted stack in Zen~5 bears on the thermal cost of the first design.
\end{itemize}

\subsection{Organisation of the paper}
\label{sec:organisation}

Part~\ref{part:foundations} reviews the architectural foundations. Part~\ref{part:chiplets} analyses chiplet-based organisation. Part~\ref{part:stacked-cache} examines 3D-stacked cache and AMD 3D V-Cache. Part~\ref{part:case-study} presents the AMD case study and quantitative analysis, and Part~\ref{part:synthesis} compares the approaches, identifies limitations and outlines future directions. Self-contained derivations longer than about half a page appear as appendices.
